\documentclass[11pt,a4paper]{article}
\usepackage[utf8]{inputenc}
\usepackage{amsmath}
\usepackage{amsfonts}
\usepackage{amssymb}
\usepackage{graphicx}
\usepackage{geometry}
\usepackage{hyperref}
\usepackage{listings}
\usepackage{color}

\geometry{left=2cm,right=2cm,top=2cm,bottom=2cm}

\title{CS4031 - Compiler Construction\\Assignment 3: Bottom-Up Parser Design \& Implementation}
\author{Hussain Waseem Syed (22i-0893) \and Mohammad Abbas (22i-1409)\\Section C}
\date{\today}

\begin{document}

\maketitle
\tableofcontents
\newpage

\section{Introduction}
Bottom-up parsing is a technique used in compiler design to construct a parse tree for an input string beginning at the leaves (the bottom) and working up towards the root (the top). Unlike top-down parsing, which attempts to find a leftmost derivation, bottom-up parsing essentially traces a rightmost derivation in reverse. LR parsers (Left-to-right, Rightmost derivation) are the most powerful class of bottom-up parsers, capable of handling a very large class of context-free grammars. In this assignment, we implement two variants: SLR(1) (Simple LR) and LR(1). SLR(1) uses global FOLLOW sets to decide when to perform reductions, which makes it simple but susceptible to conflicts in more complex grammars. LR(1) embeds specific lookahead symbols into the parser states (items), providing a more precise context for reductions, thereby resolving conflicts that SLR(1) cannot handle, albeit at the cost of a larger number of states.

\section{Approach}
\subsection{Data Structures}
\begin{itemize}
    \item \textbf{Grammar}: Modeled using \texttt{Production} objects containing a left-hand side (LHS) non-terminal and a list of right-hand side (RHS) symbols. Sets for Terminals, Non-Terminals, FIRST, and FOLLOW are maintained using HashSets for efficient O(1) lookups.
    \item \textbf{Items}: 
    \begin{itemize}
        \item \texttt{LR0Item}: Represents an SLR item. Contains a reference to a \texttt{Production} and an integer \texttt{dotIndex} tracking the progress.
        \item \texttt{LR1Item}: Composes an \texttt{LR0Item} and a \texttt{Set<String>} of lookahead symbols. 
    \end{itemize}
    \item \textbf{ItemSets}: A state in the DFA, represented by a \texttt{Set} of items (either LR0 or LR1). Assigned a unique \texttt{stateId}.
    \item \textbf{ParsingTable}: A 2D structure mapping \texttt{(State, Symbol)} to an \texttt{Action}. Actions can be \texttt{SHIFT(state)}, \texttt{REDUCE(production)}, \texttt{ACCEPT}, or \texttt{ERROR}. Conflicts are tracked if an entry already exists when a new action is proposed.
    \item \textbf{Stack}: A unified parsing stack storing \texttt{(Symbol, State)} pairs, facilitating the shift-reduce trace.
\end{itemize}

\subsection{Algorithm Implementations}
\textbf{CLOSURE}: For a given set of items $I$, we iteratively add items for any non-terminal immediately following the dot. For SLR(1), we just add \texttt{B $\rightarrow$ . $\gamma$}. For LR(1), we also compute the new lookahead by taking the FIRST set of the symbols following $B$ combined with the current item's lookahead. This is implemented via a worklist algorithm that processes until no new items are added to $I$.

\textbf{GOTO}: For an item set $I$ and symbol $X$, we collect all items in $I$ where $X$ immediately follows the dot, advance the dot, and then compute the CLOSURE of this new set. The canonical collections are built by tracking a list of visited sets and systematically expanding unvisited sets via GOTO on all grammar symbols.

\subsection{Design Decisions and Trade-offs}
\begin{itemize}
    \item \textbf{Language Choice}: Java was selected for its robust object-oriented features, built-in Collection framework (essential for Set operations), and garbage collection.
    \item \textbf{Unified Table vs Split ACTION/GOTO}: We conceptually combined them into a single \texttt{ParsingTable} object, using maps for easy lookup and formatted printing.
    \item \textbf{Immutability}: Items are treated as immutable where possible to allow safe hashing and set comparisons during DFA construction.
\end{itemize}

\subsection{Lookahead Handling in LR(1)}
In LR(1), an item is of the form \texttt{[A $\rightarrow$ $\alpha$ . $\beta$, a]}. The lookahead `a' is propagated during the CLOSURE operation. If we have \texttt{[A $\rightarrow$ $\alpha$ . B $\beta$, a]}, we add \texttt{[B $\rightarrow$ . $\gamma$, b]} for each \texttt{b $\in$ FIRST($\beta$a)}. This logic is strictly enforced in \texttt{Items.java}. Reductions in the parsing table are then \textit{only} registered for the exact lookaheads stored in the \texttt{LR1Item}, as opposed to SLR(1) which indiscriminately reduces for all symbols in \texttt{FOLLOW(A)}.

\section{Challenges and Solutions}
\begin{itemize}
    \item \textbf{Challenge}: Computing FIRST sets with $\epsilon$-productions, especially dealing with chains of nullable non-terminals.
    \item \textbf{Solution}: We implemented an iterative fixed-point algorithm that continuously passes through all productions, propagating FIRST sets and $\epsilon$-flags until no sets change size.
    \item \textbf{Challenge}: Distinguishing between shared items in LR(1) sets.
    \item \textbf{Solution}: Rather than creating a new item for every lookahead, \texttt{LR1Item} maintains a \texttt{Set} of lookaheads. When merging sets during CLOSURE, if an item with the same core exists, its lookaheads are unioned.
    \item \textbf{Challenge}: Infinite loops during CLOSURE computation due to left-recursive rules.
    \item \textbf{Solution}: Implemented equality and hash code methods properly for Items, leveraging \texttt{HashSet.contains()} to avoid reprocessing items that are already in the closure set.
\end{itemize}

\section{Test Cases}
We tested 4 distinct grammars, with multiple valid and invalid strings.
\begin{enumerate}
    \item \textbf{Grammar 1 (Simple Expressions)}: \texttt{Expr $\rightarrow$ Expr + Term | Term}, etc. (No * or parentheses).
    \begin{itemize}
        \item \textit{Valid}: \texttt{id}, \texttt{id + id}, \texttt{id + id + id} (Accepted perfectly, generated right-heavy left-associative trees).
        \item \textit{Invalid}: \texttt{id * id}, \texttt{( id )}, \texttt{+ id} (Rejected with precise syntax errors).
    \end{itemize}
    \item \textbf{Grammar 2 (Full Expressions)}: Includes multiplication and parentheses.
    \begin{itemize}
        \item \textit{Valid}: \texttt{id + id * id}, \texttt{( id + id ) * id}. Validates operator precedence.
        \item \textit{Invalid}: \texttt{id + + id}, \texttt{( id + id} (unclosed parens).
    \end{itemize}
    \item \textbf{Grammar 3 (Classic SLR Conflict)}: \texttt{Start $\rightarrow$ Lhs = Rhs | Rhs}, \texttt{Lhs $\rightarrow$ * Rhs | id}, \texttt{Rhs $\rightarrow$ Lhs}.
    \begin{itemize}
        \item \textit{Valid Strings}: \texttt{* id = id}, \texttt{id = * id}.
        \item \textit{Result}: SLR(1) generates a Shift/Reduce conflict (\texttt{s8/r5}) on \texttt{=} because \texttt{=} is in \texttt{FOLLOW(Rhs)}. LR(1) avoids this conflict because the lookahead restricts the reduction.
    \end{itemize}
    \item \textbf{Grammar 4 (Dangling Else)}: \texttt{Stmt $\rightarrow$ if Expr then Stmt else Stmt | if Expr then Stmt | other}.
    \begin{itemize}
        \item \textit{Result}: Both SLR(1) and LR(1) generate conflicts on \texttt{else}, correctly identifying the grammar as inherently ambiguous.
    \end{itemize}
\end{enumerate}

\section{Comparison Analysis}
Using \textbf{Grammar 3} as the benchmark for comparison:
\begin{table}[h]
\centering
\begin{tabular}{|l|c|c|}
\hline
\textbf{Metric} & \textbf{SLR(1)} & \textbf{LR(1)} \\ \hline
Number of states & 10 & 14 \\ \hline
ACTION entries & 17 & 22 \\ \hline
GOTO entries & 7 & 9 \\ \hline
Total table entries & 24 & 31 \\ \hline
Build time (ms) & $\sim$4 ms & $\sim$2 ms \\ \hline
Has conflicts? & \textbf{YES} (Shift/Reduce) & \textbf{No} \\ \hline
Parsing Power & Weak (Global FOLLOW) & Strong (Item Lookahead) \\ \hline
Memory Usage & Lower (fewer states) & Higher (more states) \\ \hline
\end{tabular}
\caption{SLR(1) vs LR(1) Performance Profile (Grammar 3)}
\end{table}

\textbf{Analysis}: LR(1) is strictly more powerful. In Grammar 3, SLR(1) sees the item \texttt{Rhs $\rightarrow$ Lhs .} in state 4. It looks at \texttt{FOLLOW(Rhs)}, which contains \texttt{=}, so it decides to reduce on \texttt{=}. However, an assignment cannot start with a right-hand side. LR(1) resolves this because the lookahead for that specific item is only \texttt{\$}, not \texttt{=}, preventing the erroneous reduction.

\section{Sample Outputs}
\subsection{Conflict Detection (Grammar 3 SLR trace)}
\begin{verbatim}
[String 3] tokens: [id, =, id]
=== SLR(1) ===
Step  | Stack (symbol:state)   | Remaining Input   | Action                             
-----------------------------------------------------------------------
1     | [$:0]                  | id = id $         | shift 2                            
2     | [$:0] [id:2]           | = id $            | reduce Lhs -> id                   
3     | [$:0] [Lhs:4]          | = id $            | CONFLICT: s8/r5                    

✘ String REJECTED
\end{verbatim}

\subsection{Conflict Resolution (Grammar 3 LR(1) trace)}
\begin{verbatim}
=== LR(1) ===
Step  | Stack (symbol:state)   | Remaining Input   | Action                             
-----------------------------------------------------------------------
...
3     | [$:0] [Lhs:4]          | = id $            | shift 8                            
4     | [$:0] [Lhs:4] [=:8]    | id $              | shift 10                           
...
8     | [$:0] [Start:3]        | $                 | accept                             

✔ String ACCEPTED
\end{verbatim}

\subsection{Parse Tree Generation}
Parse tree successfully generated for \texttt{( id + id ) * id} in Grammar 2:
\begin{verbatim}
Expr
└── Term
    ├── Term
    │   └── Factor
    │       ├── (
    │       ├── Expr
    │       │   ├── Expr
    │       │   │   └── Term
    │       │   │       └── Factor
    │       │   │           └── id
    │       │   ├── +
    │       │   └── Term
    │       │       └── Factor
    │       │           └── id
    │       └── )
    ├── *
    └── Factor
        └── id
\end{verbatim}

\section{Conclusion}
This assignment highlighted the practical differences in power and complexity between SLR(1) and LR(1) parsing methodologies. While SLR(1) is conceptually simpler and yields much smaller parsing tables, its reliance on global \texttt{FOLLOW} sets makes it too blunt for grammars with overlapping right contexts (like the L-value/R-value assignment grammar). LR(1) eliminates these spurious conflicts by tightly coupling contextual lookaheads to individual parser items. This precision significantly expands the class of grammars that can be parsed, albeit at the expense of computational and memory overhead to construct the larger state machine. Implementing both from scratch solidified our understanding of the shift-reduce state machine architecture.

\end{document}
